\clearpage
\section*{Anexo A. Diámetros comerciales de acero para el examen}
\addcontentsline{toc}{section}{Anexo A. Diámetros comerciales de acero para el examen}

\begin{infobox}
Para fines de este examen, el diámetro comercial se identificará con el \emph{nominal pipe size} (NPS) y su equivalente aproximado en diámetro nominal (DN). El diseño hidráulico debe realizarse siempre con \textbf{diámetro interior}, no con diámetro nominal ni con diámetro exterior.
\end{infobox}

\begin{table}[htbp]
\centering
\small
\caption{Diámetros comerciales de acero para referencia del examen}
\label{tab:nps-dn}
\begin{tabular}{ccc}
\toprule
\textbf{NPS (in)} & \textbf{DN aproximado} & \textbf{Diámetro exterior (mm)} \\
\midrule
3 & 80 & 88.9 \\
4 & 100 & 114.3 \\
5 & 125 & 141.3 \\
6 & 150 & 168.3 \\
8 & 200 & 219.1 \\
10 & 250 & 273.0 \\
12 & 300 & 323.9 \\
14 & 350 & 355.6 \\
16 & 400 & 406.4 \\
18 & 450 & 457.2 \\
20 & 500 & 508.0 \\
22 & 550 & 558.8 \\
24 & 600 & 609.6 \\
26 & 650 & 660.4 \\
28 & 700 & 711.2 \\
30 & 750 & 762.0 \\
32 & 800 & 812.8 \\
34 & 850 & 863.6 \\
36 & 900 & 914.4 \\
\bottomrule
\end{tabular}
\end{table}

\clearpage
\section*{Anexo B. Cuadro simplificado de clase de tubería de acero para el examen}
\addcontentsline{toc}{section}{Anexo B. Cuadro simplificado de clase de tubería de acero para el examen}

\begin{warnbox}
Este cuadro corresponde al \textbf{Cuadro~2, Cuadro simplificado de clase de tubería de acero para el examen}, y se provee exclusivamente para este examen. Para seleccionar la tubería, use solo los valores de esta tabla. No sustituya estos datos por catálogos externos. La presión admisible de la clase seleccionada debe ser mayor o igual que la suma de la presión estática, la sobrepresión transitoria y el margen adoptado por el diseñador.
\end{warnbox}

\begin{table}[htbp]
\centering
\footnotesize
\caption{Cuadro simplificado de clase de tubería de acero para el examen}
\label{tab:clases-acero}
\begin{tabular}{cccccc|cccccc}
\toprule
\makecell{\textbf{NPS}\\\textbf{(in)}} &
\textbf{DN} &
\textbf{Clase} &
\makecell{\textbf{$t$}\\\textbf{(mm)}} &
\makecell{\textbf{$D_i$}\\\textbf{(mm)}} &
\makecell{\textbf{$P_{\mathrm{adm}}$}\\\textbf{(MPa)}} &
\makecell{\textbf{NPS}\\\textbf{(in)}} &
\textbf{DN} &
\textbf{Clase} &
\makecell{\textbf{$t$}\\\textbf{(mm)}} &
\makecell{\textbf{$D_i$}\\\textbf{(mm)}} &
\makecell{\textbf{$P_{\mathrm{adm}}$}\\\textbf{(MPa)}} \\
\midrule
3  & 80  & 150 & 4.78  & 79.34  & 1.03 & 16 & 400 & 150 & 7.92  & 390.56 & 1.03 \\
3  & 80  & 200 & 5.56  & 77.78  & 1.38 & 16 & 400 & 200 & 9.53  & 387.34 & 1.38 \\
3  & 80  & 250 & 6.35  & 76.20  & 1.72 & 16 & 400 & 250 & 11.11 & 384.18 & 1.72 \\
4  & 100 & 150 & 4.78  & 104.74 & 1.03 & 18 & 450 & 150 & 7.92  & 441.36 & 1.03 \\
4  & 100 & 200 & 5.56  & 103.18 & 1.38 & 18 & 450 & 200 & 9.53  & 438.14 & 1.38 \\
4  & 100 & 250 & 6.35  & 101.60 & 1.72 & 18 & 450 & 250 & 11.11 & 434.98 & 1.72 \\
6  & 150 & 150 & 4.78  & 158.74 & 1.03 & 20 & 500 & 150 & 9.53  & 488.94 & 1.03 \\
6  & 150 & 200 & 5.56  & 157.18 & 1.38 & 20 & 500 & 200 & 11.11 & 485.78 & 1.38 \\
6  & 150 & 250 & 6.35  & 155.60 & 1.72 & 20 & 500 & 250 & 12.70 & 482.60 & 1.72 \\
8  & 200 & 150 & 6.35  & 206.40 & 1.03 & 24 & 600 & 150 & 9.53  & 590.54 & 1.03 \\
8  & 200 & 200 & 7.92  & 203.26 & 1.38 & 24 & 600 & 200 & 11.11 & 587.38 & 1.38 \\
8  & 200 & 250 & 9.53  & 200.04 & 1.72 & 24 & 600 & 250 & 12.70 & 584.20 & 1.72 \\
10 & 250 & 150 & 6.35  & 260.30 & 1.03 & 30 & 750 & 150 & 11.11 & 739.78 & 1.03 \\
10 & 250 & 200 & 7.92  & 257.16 & 1.38 & 30 & 750 & 200 & 12.70 & 736.60 & 1.38 \\
10 & 250 & 250 & 9.53  & 253.94 & 1.72 & 30 & 750 & 250 & 14.27 & 733.46 & 1.72 \\
12 & 300 & 150 & 6.35  & 311.20 & 1.03 & 36 & 900 & 150 & 12.70 & 889.00 & 1.03 \\
12 & 300 & 200 & 7.92  & 308.06 & 1.38 & 36 & 900 & 200 & 14.27 & 885.86 & 1.38 \\
12 & 300 & 250 & 9.53  & 304.84 & 1.72 & 36 & 900 & 250 & 15.88 & 882.64 & 1.72 \\
14 & 350 & 150 & 7.92  & 339.76 & 1.03 &    &     &     &       &        &      \\
14 & 350 & 200 & 9.53  & 336.54 & 1.38 &    &     &     &       &        &      \\
14 & 350 & 250 & 11.11 & 333.38 & 1.72 &    &     &     &       &        &      \\
\bottomrule
\end{tabular}
\caption*{\footnotesize Donde $t$ es el espesor mínimo, $D_i$ es el diámetro interior de diseño y $P_{\mathrm{adm}}$ es la presión admisible de la clase.}
\end{table}

\clearpage
\section*{Anexo C. Presión admisible del suelo traducida a SI}
\addcontentsline{toc}{section}{Anexo C. Presión admisible del suelo traducida a SI}

\begin{infobox}
Las presiones admisibles de suelo de referencia para el diseño de bloques de anclaje se presentan en el \textbf{Cuadro~3}. Para este examen, use como caso base la fila de \textbf{mezcla densa de arena y grava} a \SI{1}{\meter} de profundidad, con $\phi = 40^\circ$ y $\sigma_a = \SI{0.0862}{\mega\pascal}$.
\end{infobox}

\begin{table}[htbp]
\centering
\footnotesize
\caption{Presión admisible del suelo traducida a unidades SI}
\label{tab:presion-suelo}
\begin{tabular}{lcccc}
\toprule
\textbf{Material natural del suelo} & \textbf{0.61 m} & \textbf{0.91 m} & \textbf{1.22 m} & \textbf{1.52 m} \\
 & \multicolumn{4}{c}{\textbf{Presión admisible (kPa)}} \\
\midrule
Roca sana & 383.1 & 478.8 & 478.8 & 478.8 \\
Mezcla densa de arena y grava ($\phi = 40^\circ$) & 57.4 & 86.2 & 114.9 & 143.6 \\
Arena densa fina a gruesa ($\phi = 35^\circ$) & 38.3 & 57.4 & 79.0 & 100.6 \\
Mezcla de limo y arcilla ($\phi = 25^\circ$) & 23.9 & 33.5 & 45.5 & 57.4 \\
Arcilla blanda y suelos orgánicos ($\phi = 10^\circ$) & 9.6 & 14.4 & 19.2 & 23.9 \\
\bottomrule
\end{tabular}
\end{table}

\clearpage
\section*{Anexo D. Valores de k para análisis de golpe de ariete}
\addcontentsline{toc}{section}{Anexo D. Valores de k para análisis de golpe de ariete}

\begin{table}[htbp]
\centering
\footnotesize
\caption{Valores de $k$ tomados de mecánica de sólidos para el análisis de golpe de ariete}
\label{tab:k-golpe-ariete}
\begin{tabularx}{\textwidth}{>{\raggedright\arraybackslash}p{0.38\textwidth} >{\centering\arraybackslash}X >{\centering\arraybackslash}p{0.18\textwidth}}
\toprule
\textbf{Condiciones de apoyo} & \textbf{$k$ (paredes gruesas)} & \textbf{$k$ (paredes delgadas)} \\
\midrule
Anclada solo en el extremo aguas arriba. Libre para moverse longitudinalmente &
$\dfrac{2t}{D}(1+\mu)+\left(\dfrac{D}{D+t}\right)-\left(1-\dfrac{\mu}{2}\right)$ &
$1-\dfrac{\mu}{2}$ \\
Anclada contra movimiento axial (longitudinal), de modo que el esfuerzo axial sea cero &
$\dfrac{2t}{D}(1+\mu)+\left(\dfrac{D}{D+t}\right)(1-\mu^2)$ &
$1-\mu^2$ \\
Con juntas de expansión en toda la longitud (esfuerzo axial = 0) &
$\dfrac{2t}{D}(1+\mu)+\left(\dfrac{D}{D+t}\right)(1)$ &
$1$ \\
En túneles ($t=\infty$) &
$\dfrac{2t}{D}(1+\mu)$ &
No aplica \\
\bottomrule
\end{tabularx}
\caption*{\footnotesize Nota: $\dfrac{D}{D+t}= \dfrac{1}{1+t/D}$; $\mu$ es la razón de Poisson y típicamente toma valores entre 0.25 y 0.30 para muchos materiales de tubería.}
\end{table}

\clearpage
\section*{Anexo E. Coeficientes de pérdida por codos según Hinds/USBR}
\addcontentsline{toc}{section}{Anexo E. Coeficientes de pérdida por codos según Hinds/USBR}

\begin{figure}[h]
\centering
\rotatebox{90}{\includegraphics[page=1,scale=0.73, trim={2.5cm 1.8cm 2.5cm 2cm},clip]{../insumos/Head loss coefficients for pipe bends.pdf}}
\caption{Anexo de referencia para coeficientes de pérdida por codos según Hinds/USBR, citado en la sección de pérdidas menores y accesorios.}
\label{fig:hinds-usbr}
\end{figure}

\clearpage
\section*{Anexo F. Encabezados mínimos de la tabla entregable}
\addcontentsline{toc}{section}{Anexo F. Encabezados mínimos de la tabla entregable}

\begin{longtable}{p{0.95\textwidth}}
\toprule
\textbf{Encabezados sugeridos para la tabla principal del entregable} \\
\midrule
Ítem \\
Progresiva inicial \\
Progresiva final \\
Longitud del tramo (\si{\meter}) \\
Tipo de elemento \\
Diámetro nominal \\
Diámetro interior de diseño (\si{\milli\meter}) \\
Material \\
Clase de tubería \\
Caudal (\si{\liter\per\second} o \si{\meter\cubed\per\second}) \\
Velocidad (\si{\meter\per\second}) \\
Pérdida por fricción (\si{\meter}) \\
Pérdida menor (\si{\meter}) \\
Pérdida total acumulada (\si{\meter}) \\
Cota del terreno (\si{\meter}) \\
Cota de la tubería (\si{\meter}) \\
Carga piezométrica (\si{\meter}) \\
Carga total o línea de energía (\si{\meter}) \\
Presión (\si{\meter} de columna de agua o \si{\kilo\pascal}) \\
Pieza especial / válvula / ventosa / purga \\
Observaciones \\
\bottomrule
\end{longtable}
